\scanpage{8}
a aussi l'inégalité suivante (résultant déjà de la suite exacte des
trois derniers termes de (8.5.1)):
\[
(8.5.6)\qquad \dim(E(Y)\otimes_W k)_{\mathrm{SS}}\geq\dim E(Y,\underline O_Y)_{\mathrm{SS}}\deleted{+\dim\operatorname{Ker}w}
\]
\deleted{\[\dim E(Y,\underline O_Y)_{\mathrm{SS}}-t^{n-1}(X,p),\]
où $t^i(X,p)$ désigne le nombre de \inserted{$p$-}coefficients de torsion de $X$ (défini
via la cohomologie cristalline de $X$ en dimension $i$).}

\uline{Remarques} 8.5.7. Je pense que l'homomorphisme $v$ de (8.5.2) est
injectif\deleted{\unsure{,}} (ce qui simplifie en conséquence (8.5.1) et (8.5.5), puisque
on aurait $\operatorname{Ker}v=0$); c'est vrai en tous cas si \inserted{$n=2$, ou si}
$Y$ est \deleted{de degré assez}
une section par une hypersurface de degré $d$ assez grand pour qu'on ait
\[H^{n-1}(X,\underline O_X(-d))=0.\]

\deleted{D'autre part, on peut si on veut dans 8.5.5 bloquer ensemble les \unsure{deux}
termes $\dim\operatorname{Coker}w-\dim\operatorname{Ker}w$, \deleted{car on aura comme}\inserted{pour trouver} la différence des
rangs \deleted{$\dim\operatorname{Coker}w-\dim\operatorname{Ker}w$} des deux membres de (8.5.3).}

8.5.8. Lorsque \deleted{dans} l'on applique (8.5.5) en prenant pour $Y$ une secti-
on de $X$ avec une hypersurface de degré $d$ ``assez générale'' (de telle
façon que le rang semi-simple de $E(Y,\underline O_Y)$ ait la valeur maxima), alors
on trouve qu'aux terme $-\dim\operatorname{Ker}v$ près (sans doute toujours nul, cf.
8.5.\unsure{7}) le deuxième membre est \uline{indépendant du choix} \deleted{de} \uline{particulier de $Y$}.
Cette remarque s'applique en particulier (8.5.7) si $d$ est assez grand.
On en conclut que si $k$ est fini, \deleted{\illegible{1}} on peut trouver un ouvert \inserted{$U_d$} de l'es\unsure{pa}-
ce projectif des hypersurfaces de degré $d$, tel que pour \deleted{tout}\inserted{un} point \inserted{\unsure{variable}} de $U$
à coefficients dans une extension finie $k'$ de $k$, le nombre de valeurs
propres de Frobenius opérant sur la cohomologie évanescente de $Y_{\bar k}$ et
qui sont des unités $p$-adiques soit constant. \deleted{Mieux, quitte à agrandir
un peu $k$, de sorte que les valeurs propres de $\operatorname{Frob}_k$} Nous désirons mon-
\mspage{7}
